\documentclass[11pt]{article}
\usepackage[margin=1in]{geometry}
\usepackage{amsmath,amssymb,amsthm}
\usepackage[colorlinks=true,urlcolor=blue]{hyperref}
\newtheorem{theorem}{Theorem}
\title{A counterexample to the \(R\)-difference\\infinite log-concavity conjecture}
\author{Independent working note --- MathDB 331961}
\date{12 September 2026 (UTC)}
\begin{document}
\maketitle
\paragraph{Task 1: resolve the conjecture.}
Let
\[
 R_n=\sum_{k=0}^n\binom nk\binom{n+k}k\frac1{2k-1},
 \qquad a_n=R_n^2-R_{n-1}R_{n+1}\quad(n\ge6).
\]
Conjecture 4.8 in the published paper \cite{sun} asserts that
\(\{a_n\}_{n\ge6}\) is infinitely log-concave.
For the log-concavity operator
\[
 (\mathcal L a)_n=a_n^2-a_{n-1}a_{n+1},
\]
we give an exact negative value at the first available center.

\begin{theorem}
The published Conjecture 4.8 is false. In fact,
\[
 (\mathcal L a)_7=-343413129984<0.
\]
Both \(\{-a_n\}_{n\ge6}\) and \(\{|a_n|\}_{n\ge6}\) fail the same
log-concavity test.
\end{theorem}
\begin{proof}
Direct evaluation of the defining finite binomial sums gives
\[
\begin{array}{c|rrrrr}
 n&5&6&7&8&9\\ \hline
 R_n&329&1359&6001&27759&132689.
\end{array}
\]
Consequently
\begin{align*}
 a_6&=1359^2-329\cdot6001=-127448,\\
 a_7&=6001^2-1359\cdot27759=-1712480,\\
 a_8&=27759^2-6001\cdot132689=-25704608.
\end{align*}
All three indices are within the stated subsequence. Thus
\begin{align*}
 a_7^2-a_6a_8
 &=1712480^2-127448\cdot25704608\\
 &=-343413129984<0.
\end{align*}
The sequence is not even log-concave, so it cannot be infinitely
log-concave. Reversing every sign does not change this quadratic
expression. Taking absolute values likewise leaves it unchanged
because the three displayed terms have the same sign.
\end{proof}

\paragraph{Iteration convention.}
Even if one tests log-concavity only after first applying \(\mathcal L\),
the assertion fails. Writing \(b=\mathcal L a\), direct calculation gives
\[
 b_7=-343413129984,\quad b_8=-78797107789056,\quad
 b_9=-18331042187598528,
\]
and hence
\[
 b_8^2-b_7b_9=-86136377591851264277692416<0.
\]
Only terms of \(a\) with indices at least \(6\) occur. These
calculations also use \(R_{10}=649815\) and \(R_{11}=3242377\).

\paragraph{Verification and scope.}
The standard-library script \texttt{verify\_counterexample.py} evaluates
the binomial sums with rational arithmetic and independently checks
them against the source recurrence
\[
 (n+3)R_{n+3}-(7n+13)R_{n+2}+(7n+15)R_{n+1}-(n+1)R_n=0,
\]
starting from \(R_0=-1,R_1=1,R_2=7\).
The complete exact certificate includes the individual summands.
The counterexample matches both the current arXiv source and
the published Conjecture 4.8. It does not depend on a negative-term
convention, a reversed sign, or an endpoint outside the sequence.
The dated source and public-claim audit is saved separately.

\begin{thebibliography}{9}
\bibitem{sun} Brian Yi Sun,
\emph{On a ratio monotonicity conjecture of a new kind of numbers},
Journal of Inequalities and Applications \textbf{2018}, article 24,
Conjecture 4.8.
\href{https://doi.org/10.1186/s13660-018-1614-1}{doi:10.1186/s13660-018-1614-1}.
Earlier version: \href{https://arxiv.org/abs/1512.01008}{arXiv:1512.01008}.
\end{thebibliography}

\newpage
\paragraph{Task 2: determine the actual log-convexity range.}
Set \(d_n=R_{n-1}R_{n+1}-R_n^2\).
Determine the least \(N\) for which \(\{d_n\}_{n\ge N}\) is positive
and strictly log-convex. Give a proof covering every subsequent index.

\paragraph{Task 3: test infinite log-convexity as a replacement.}
For \((\mathcal C d)_n=d_{n-1}d_{n+1}-d_n^2\), determine whether
there is one finite starting index such that every iterate
\(\mathcal C^j d\) is positive at all of its remaining indices.
If each order instead requires its own starting index, determine
how the least such index grows with \(j\).

\paragraph{Task 4: investigate higher Hankel determinants.}
For \(m\ge2\), determine the signs and eventual nonvanishing of
\[
 \det\bigl(R_{n+i+j}\bigr)_{0\le i,j<m}.
\]
Obtain effective thresholds in terms of \(m\), and determine whether
any fixed shift of the \(R\)-sequence has all Hankel minors nonnegative.

\paragraph{Task 5: classify normalizations of the difference.}
Determine the real parameters \(\alpha\) for which the positive
sequence \(\{n^\alpha|a_n|\}_{n\ge6}\) is log-concave.
Distinguish the exact range for this fixed starting index from the
range giving eventual log-concavity.
\end{document}
